\begin{figure}[h]
\begin{inductive}{\Q \vdash \ell \sim \tau}
  \inferrule
    { \mathrm{bits}(T) = s \\ T\ \text{is a primitive integer or boolean type} }
    { \Q \vdash \langle o, s \rangle \sim T }
    {\textsc{PrimMatch}}

  \quad
  \inferrule
    { \text{for each } i,\ \Q \vdash \ell_i \sim \tau_i }
    { \Q \vdash \texttt{record}\ \{\, f_i : \ell_i \,\} \sim
        \texttt{struct}\ \{\, f_i : \tau_i \,\} }
    {\textsc{RecordMatch}}

  \quad
  \inferrule
    { \Q \vdash \ell \sim \tau }
    { \Q \vdash \ell [n] \sim \tau [n] }
    {\textsc{ArrayMatch}}

  \quad
  \inferrule
    { }
    { \Q \vdash () \sim () }
    {\textsc{UnitMatch}}

\end{inductive}

\[
\begin{array}{ll}
  \mathrm{bits}(U8) = 8 \quad & \mathrm{bits}(U16) = 16 \\
  \mathrm{bits}(U32) = 32 \quad & \mathrm{bits}(U64) = 64 \\
  \mathrm{bits}(\_Bool) = 1
\end{array}
\]
\caption{Layout–type matching rules in Cedar. Each layout $\ell$ corresponds to a well-typed region of memory for type $\tau$. \label{fig:cedar-match}}
\end{figure}

Figure~\ref{fig:cedar-match} summarizes the core typing rules.
This project implements a subset of Dargent's layout-type correspondence
that applies to fixed-type size, byte addressed C records and arrays.